% !TeX root = ../../main.tex
% =====================================================================
%  Proceso de diseño y toma de decisiones
%  Responsable: 
%  Última revisión: 
% =====================================================================
\chapter{Proceso de diseño y toma de decisiones}
\label{ch:proceso_diseno}

\begin{objetivos}
  \item \redactar{Objetivo de aprendizaje 1 --- concreto y comprobable.}
  \item \redactar{Objetivo 2.}
  \item \redactar{Objetivo 3.}
\end{objetivos}

\fichasesion{90 min}{\redactar{capítulos previos}}{\redactar{material}}

\section{Fases del ciclo de diseño}
\label{sec:proceso_diseno:01}
\pendiente{Texto del apartado pendiente de redactar. Las figuras
\ref{fig:proceso_diseno} y \ref{fig:subproceso_dimensionamiento} ya están
puestas; falta el texto que las explica.}

La figura \ref{fig:proceso_diseno} recoge el proceso que sigue el equipo para
diseñar un componente, desde que se asigna la tarea hasta que el resultado se
entrega al encargado de área. En trazo continuo está el proceso tal y como se
planteó; en trazo discontinuo rojo, los pasos propuestos que todavía no se han
acordado y que se justifican al final del apartado.

Los triángulos numerados son enlaces a subprocesos que se desarrollan por
separado. El único desarrollado hasta ahora es el \textsc{ii}, el
dimensionamiento preliminar, que recoge la figura
\ref{fig:subproceso_dimensionamiento}.

% ---------------------------------------------------------------------
\begin{figure}[p]
\centering
\resizebox{!}{0.90\textheight}{%
\begin{tikzpicture}[
    proceso/.append style={font=\sffamily\footnotesize},
    decision/.append style={font=\sffamily\footnotesize},
    terminal/.append style={font=\sffamily\footnotesize\bfseries,text width=44mm},
    nota/.append style={font=\sffamily\scriptsize},
    conector/.append style={font=\sffamily\scriptsize\bfseries,minimum size=9mm},
    etiq/.append style={font=\sffamily\scriptsize},
    cabecera/.style={proceso,text width=24mm}]

% ---- Flujo principal ----
\node[terminal] (t1)  at (0,0)      {Asignación de tareas};

\node[cabecera] (b1) at (-4.65,-1.75) {Documentos aplicables};
\node[cabecera] (b2) at (-1.55,-1.75) {Objetivos};
\node[cabecera] (b3) at ( 1.55,-1.75) {Método de cálculo};
\node[cabecera,propuesto] (b4) at ( 4.65,-1.75) {Interfaces y restricciones};

\node[proceso]  (p1) at (0,-3.55)  {Dimensionamiento preliminar};
\node[conector] (c2) at (2.7,-3.55) {II};
\node[nota]     (n1) at (6.9,-3.55)
      {\textbf{Pasos}\\[1pt]$\cdot$ Selección de material\\
       $\cdot$ Método de fabricación\\$\cdot$ Ubicación de componentes};

\node[decision] (d1) at (0,-5.35)  {¿Resultado válido?};
\node[proceso,propuesto] (p2) at (0,-7.05) {Registro de decisión};

\node[proceso]  (p3) at (0,-8.75)  {Diseño CAD};
\node[conector] (c3) at (2.7,-8.75) {III};
\node[nota]     (n2) at (6.9,-8.75)
      {\textbf{Pasos}\\[1pt]$\cdot$ Ensamblaje\\$\cdot$ Fabricación\\
       $\cdot$ Material\\$\cdot$ Normativa};

\node[decision] (d2) at (0,-10.55) {¿CAD válido?};

\node[proceso]  (p4) at (0,-12.25) {Simulación MEF};
\node[conector] (c4) at (2.7,-12.25) {IV};
\node[nota]     (n3) at (6.9,-12.25)
      {\textbf{Pasos}\\[1pt]$\cdot$ Condiciones de contorno\\
       $\cdot$ Tipo de vinculación\\$\cdot$ Mallado\\$\cdot$ Convergencia};

\node[decision] (d3) at (0,-14.05) {¿Simulación válida?};
\node[proceso]  (p5) at (0,-15.85) {Revisión de objetivos};
\node[decision] (d4) at (0,-17.65) {¿Objetivos cumplidos?};
\node[nota]     (n4) at (-5.2,-17.65)
      {Revisión completa\\Detección de fallos\\Iteración};

\node[proceso,text width=42mm] (p6) at (0,-19.75)
      {Elaboración de BOM $\cdot$ Informe de diseño $\cdot$ Informe de coste
       $\cdot$ Informe de simulaciones $\cdot$ Hoja de procesos};
\node[conector] (c5) at (4.4,-19.75) {V};

\node[decision,propuesto] (d5) at (0,-21.85) {¿Revisado por otra persona?};
\node[terminal] (t2) at (0,-23.65) {Enviar al encargado de área};
\node[proceso,propuesto,text width=34mm] (p7) at (0,-25.35)
      {Fabricación y validación experimental};

% ---- Flechas del tronco ----
\draw[flujo] (t1.south) -- ++(0,-0.35) coordinate (r1);
\draw[draw=fusalPrim!85,semithick] (b1.north|-r1) -- (b4.north|-r1);
\foreach \n in {b1,b2,b3} \draw[flujo] (\n.north|-r1) -- (\n.north);
\draw[flujoprop] (b4.north|-r1) -- (b4.north);

\coordinate (r2) at (0,-2.65);
\foreach \n in {b1,b2,b3} \draw[draw=fusalPrim!85,semithick] (\n.south) -- (\n.south|-r2);
\draw[draw=fusalSec,dashed,dash pattern=on 2pt off 1.5pt] (b4.south) -- (b4.south|-r2);
\draw[draw=fusalPrim!85,semithick] (b1.south|-r2) -- (b4.south|-r2);
\draw[flujo] (r2) -- (p1.north);

\draw[flujo] (p1) -- (d1);
\draw[flujoprop] (d1) -- node[etiq,pos=0.45]{Sí} (p2);
\draw[flujoprop] (p2) -- (p3);
\draw[flujo] (d2) -- node[etiq,pos=0.45]{Sí} (p4);
\draw[flujo] (p3) -- (d2);
\draw[flujo] (d3) -- node[etiq,pos=0.45]{Sí} (p5);
\draw[flujo] (p4) -- (d3);
\draw[flujo] (p5) -- (d4);
\draw[flujo] (d4) -- node[etiq,pos=0.40]{Sí} (p6);
\draw[flujoprop] (p6) -- (d5);
\draw[flujoprop] (d5) -- node[etiq,pos=0.45]{Sí} (t2);
\draw[flujoprop] (t2) -- (p7);

% ---- Conectores y notas ----
\foreach \p/\c in {p1/c2,p3/c3,p4/c4} \draw[aviso] (\p.east) -- (\c.west);
\draw[aviso] (p6.east) -- (c5.west);
\foreach \c/\n in {c2/n1,c3/n2,c4/n3} \draw[aviso] (\c.east) -- (\n.west);
\draw[aviso] (d4.west) -- node[etiq,pos=0.45]{No} (n4.east);

% ---- Realimentaciones ----
\draw[flujo] (d1.west) -- node[etiq,pos=0.3]{No} ++(-2.1,0) |- (p1.west);
\draw[flujo] (d2.west) -- node[etiq,pos=0.3]{No} ++(-2.6,0) |- ([yshift=-1.2mm]p1.west);
\draw[flujo] (d3.west) -- node[etiq,pos=0.3]{No} ++(-3.1,0) |- ([yshift=1.2mm]p1.west);
\draw[flujoprop] (n4.west) -- ++(-1.1,0) |- ([yshift=-2.4mm]p1.west);
\draw[flujoprop] (d5.east) -- node[etiq,pos=0.35]{No} ++(2.0,0) |- (p6.east);
\draw[flujoprop] (p7.east) -- ++(3.9,0) |- (p5.east);

\end{tikzpicture}}
\caption[Proceso de diseño de un componente]{Proceso de diseño de un
componente. Trazo continuo: el proceso tal y como se planteó. Trazo discontinuo
rojo: pasos propuestos, pendientes de acordar. Los triángulos enlazan con
subprocesos desarrollados aparte.}
\label{fig:proceso_diseno}
\end{figure}

% ---------------------------------------------------------------------
\begin{figure}[p]
\centering
\begin{tikzpicture}[
    proceso/.append style={font=\sffamily\footnotesize},
    decision/.append style={font=\sffamily\footnotesize},
    terminal/.append style={font=\sffamily\footnotesize\bfseries,text width=40mm},
    nota/.append style={font=\sffamily\scriptsize},
    conector/.append style={font=\sffamily\scriptsize\bfseries,minimum size=9mm},
    etiq/.append style={font=\sffamily\scriptsize}]

\node[conector] (ci) at (0,0.9) {II};
\node[proceso]  (s1) at (0,-0.5)  {Búsqueda de referencias};
\node[nota]     (m1) at (4.6,-0.5) {$\cdot$ Libros\\$\cdot$ TFG y TFM\\
       $\cdot$ Foros\\$\cdot$ Vídeos y documentación de otros equipos};

\node[decision] (e1) at (0,-2.4)  {¿Método fiable?};
\node[proceso]  (s2) at (0,-4.3)  {Desarrollo analítico};
\node[proceso]  (s3) at (0,-6.2)  {Programar el cálculo};
\node[nota]     (m2) at (-4.6,-6.2) {$\cdot$ Excel\\$\cdot$ MathCAD\\
       $\cdot$ MATLAB\\$\cdot$ Programación};

\node[proceso]  (s4) at (0,-8.1)  {Validación y comparación};
\node[nota]     (m3) at (4.6,-8.1) {$\cdot$ Objetivos\\
       $\cdot$ Referencias anteriores\\$\cdot$ Ensayo, cuando exista};

\node[decision] (e2) at (0,-10.0) {¿Válido?};
\node[proceso,propuesto,text width=34mm] (s5) at (0,-11.9)
      {Documentar hipótesis y archivar el modelo};
\node[terminal] (fin) at (0,-13.7) {Continuar: diseño CAD};

\draw[flujo] (ci) -- (s1);
\draw[flujo] (s1) -- (e1);
\draw[flujo] (e1) -- node[etiq,pos=0.45]{Sí} (s2);
\draw[flujo] (s2) -- (s3);
\draw[flujo] (s3) -- (s4);
\draw[flujo] (s4) -- (e2);
\draw[flujoprop] (e2) -- node[etiq,pos=0.45]{Sí} (s5);
\draw[flujoprop] (s5) -- (fin);

\draw[flujo] (e1.east) -- node[etiq,pos=0.35]{No} ++(2.2,0) |- (s3.east);
\draw[flujo] (e2.west) -- node[etiq,pos=0.25]{No} ++(-5.4,0) |- (s1.west);

\draw[aviso] (s1.east) -- (m1.west);
\draw[aviso] (s3.west) -- (m2.east);
\draw[aviso] (s4.east) -- (m3.west);

\end{tikzpicture}
\caption[Subproceso \textsc{ii}: dimensionamiento preliminar]{Subproceso
\textsc{ii}: dimensionamiento preliminar. Los subprocesos \textsc{iii} (diseño
CAD), \textsc{iv} (simulación MEF) y \textsc{v} (documentación) están pendientes
de desarrollar.}
\label{fig:subproceso_dimensionamiento}
\end{figure}

\subsection{Qué falta en el proceso}

Al revisar el diagrama aparecen seis huecos. Cuatro de ellos ---los números
\ref{item:huecos:revision}, \ref{item:huecos:ddr}, \ref{item:huecos:validacion} y
\ref{item:huecos:interfaces}--- están dibujados en rojo en la figura
\ref{fig:proceso_diseno}. Los otros dos no, porque antes hay que decidir dónde
encajan.

\begin{enumerate}
  \item \textbf{No hay salida de los bucles.} Las tres preguntas de validez
        devuelven al dimensionamiento preliminar sin límite. Hace falta un
        criterio de escalado: tras dos iteraciones sin converger, el problema
        pasa al encargado de área. Si no, una tarea puede consumir media
        temporada sin que nadie se entere.
  \item\label{item:huecos:revision} \textbf{Nadie revisa el trabajo antes de
        entregarlo.} El capítulo
        \ref{ch:herramientas_cad_datos} exige que una pieza no se libere sin que
        la revise alguien distinto del autor. Esa comprobación tiene que estar en
        el proceso, no en la buena voluntad de cada uno.
  \item\label{item:huecos:ddr} \textbf{El registro de decisión no aparece.} El diagrama recoge qué se
        hace, pero no deja rastro de \emph{por qué} se eligió una alternativa y
        no otra. Sin eso, el año que viene se vuelve a empezar --- que es
        justamente lo que este libro intenta evitar.
  \item \textbf{El coste entra demasiado tarde.} Aparece solo en la
        documentación final, cuando ya no se puede cambiar nada. Debería ser un
        criterio de la primera pregunta de validez, junto con la fabricabilidad.
  \item\label{item:huecos:validacion} \textbf{No hay realimentación desde la fabricación ni desde la pista.}
        El proceso termina al entregar los documentos. Lo que pasa después ---
        que la pieza no se pueda fabricar como estaba prevista, o que rompa ---
        no vuelve a entrar en ningún sitio.
  \item\label{item:huecos:interfaces} \textbf{Las interfaces con otros subsistemas no se acuerdan.} El bloque
        de partida recoge documentos, objetivos y método, pero no con quién hay
        que hablar antes de fijar una cota compartida (capítulo
        \ref{ch:coche_como_sistema}).
\end{enumerate}

\pendiente{Decidir con el equipo cuáles de estos seis pasos se incorporan, y
actualizar la figura. Lo que se acuerde deja de ser trazo discontinuo.}

\pendiente{Transcripción de la versión manuscrita: revisar las etiquetas
«Desarrollo analítico» y «Validación y comparación», y el contenido del recuadro
de herramientas de cálculo, que se han leído con dudas.}

\pendiente{El boceto original incluye un recuadro hexagonal, enlazado al conector
\textsc{i}, que dice «Aclaración: Coordinador, Encargado y Supervisor». No se ha
llevado a la figura porque no está claro a qué paso del proceso se refiere.
Decidir si es la definición de los tres papeles que intervienen ---y entonces va
en el apartado \ref{sec:proceso_diseno:04}, junto al registro de decisiones--- o
si marca los puntos del flujo en los que cada uno interviene, en cuyo caso hay
que dibujarlo.}

\section{Generación de conceptos}
\label{sec:proceso_diseno:02}
\pendiente{Pendiente de redactar.}

\section{Matrices de decisión}
\label{sec:proceso_diseno:03}
\pendiente{Pendiente de redactar.}

\section{Registro de decisiones}
\label{sec:proceso_diseno:04}
\pendiente{Pendiente de redactar.}

\section{Congelación del diseño}
\label{sec:proceso_diseno:05}
\pendiente{Pendiente de redactar.}

% ---------------------------------------------------------------------
%  Cierre del capítulo: los cuatro recuadros son obligatorios.
% ---------------------------------------------------------------------

\begin{caso}
\redactar{Qué hicimos nosotros: decisión tomada, datos de partida, resultado en pista y qué haríamos distinto.}
\end{caso}

\begin{ojo}
\begin{itemize}
  \item \redactar{Error que repetimos cada temporada en este tema.}
\end{itemize}
\end{ojo}

\begin{entregable}
\redactar{Qué produce el lector al terminar: plantilla, hoja de cálculo, checklist o apartado del informe.}
\end{entregable}

\begin{juez}
\begin{enumerate}
  \item \redactar{Pregunta tipo del design event sobre este capítulo.}
\end{enumerate}
\end{juez}
